\documentclass[12pt]{article}
% Shared preamble and text for the data-request letters.
% Replace the bracketed fields before submitting any letter.

\usepackage[a4paper,top=2.4cm,bottom=2.4cm,left=2.8cm,right=2.4cm]{geometry}
\usepackage[utf8]{inputenc}
\usepackage[T1]{fontenc}
\usepackage{lmodern}
\usepackage{microtype}
\usepackage{setspace}
\usepackage{array}
\usepackage{booktabs}
\usepackage{enumitem}
\usepackage{tabularx}
\usepackage{xcolor}
\usepackage[hidelinks]{hyperref}

\setstretch{1.12}
\setlength{\parindent}{0pt}
\setlength{\parskip}{6pt}
\setlist[itemize]{leftmargin=1.5em,itemsep=2pt,topsep=2pt}
\setlist[enumerate]{leftmargin=1.7em,itemsep=2pt,topsep=2pt}
\newcolumntype{Y}{>{\raggedright\arraybackslash}X}
\newcolumntype{P}[1]{>{\raggedright\arraybackslash}p{#1}}

\newcommand{\NamaMahasiswa}{[NAMA LENGKAP MAHASISWA]}
\newcommand{\NIMMahasiswa}{[NIM]}
\newcommand{\NamaPembimbing}{[NAMA DAN GELAR DOSEN PEMBIMBING]}
\newcommand{\JudulSementara}{Diagnosis \textit{Borrowed Size} dan \textit{Agglomeration Shadow} antara Jakarta dan Pusat-pusat Sekunder di Kawasan Aglomerasi Jakarta}

\newcommand{\KopSurat}{%
  \begin{center}
    \fbox{%
      \begin{minipage}{0.88\textwidth}
        \centering
        \textbf{KOP SURAT RESMI FAKULTAS TEKNIK UNIVERSITAS GADJAH MADA}\\
        \small Ganti blok ini dengan kop resmi sebelum surat dikirim.
      \end{minipage}}
  \end{center}
  \vspace{4pt}
}

\newcommand{\IdentitasMahasiswa}{%
  \begin{tabularx}{\textwidth}{@{}P{0.30\textwidth}Y@{}}
    Nama & \NamaMahasiswa \\
    NIM & \NIMMahasiswa \\
    Program studi & Sarjana Perencanaan Wilayah dan Kota \\
    Departemen dan fakultas & Departemen Teknik Arsitektur dan Perencanaan, Fakultas Teknik, Universitas Gadjah Mada \\
    Dosen pembimbing & \NamaPembimbing
  \end{tabularx}
}

\newcommand{\PernyataanPenggunaanData}{%
Data yang dimohon akan digunakan untuk kepentingan pendidikan dan penelitian nonkomersial. Data tidak akan dijual, dialihkan, atau disebarluaskan dalam bentuk mentah. Hasil yang dipublikasikan akan berupa statistik agregat, peta turunan, atau keluaran lain yang diizinkan oleh pemilik data. Pengolahan dan penyimpanan data akan mengikuti ketentuan kerahasiaan, lisensi, pembatasan publikasi, dan perjanjian penggunaan data yang berlaku.
}

\newcommand{\PernyataanTanpaBiaya}{%
Apabila permintaan ini tidak dapat dipenuhi tanpa biaya, mohon kami diberi tahu terlebih dahulu. Kami tidak menyetujui pembuatan tabulasi atau pengadaan data berbayar tanpa persetujuan baru dari mahasiswa dan dosen pembimbing.
}

\newcommand{\AbstraksiPenelitian}{%
\textbf{Latar belakang.} Kegiatan ekonomi, perjalanan, layanan, dan perumahan di sekitar Jakarta berlangsung melampaui batas administrasi DKI Jakarta. Hubungan tersebut tidak otomatis menunjukkan bahwa pusat-pusat sekunder telah memperkuat fungsi lokalnya. Suatu wilayah dapat memperoleh akses ke pasar metropolitan, tetapi pada saat yang sama semakin bergantung pada Jakarta untuk pekerjaan, layanan, atau fungsi bernilai tinggi.

\textbf{Tujuan.} Penelitian ini mengkaji apakah integrasi dengan Jakarta berjalan bersama penguatan fungsi lokal, suatu pola yang konsisten dengan \textit{borrowed size}, atau bersama defisit fungsi dan ketergantungan yang asimetris, suatu pola yang konsisten dengan \textit{agglomeration shadow}. Penelitian tidak bertujuan mengusulkan perluasan wilayah DKI Jakarta, mengubah batas pemerintahan, atau membuktikan hubungan kausal hanya dari potret satu periode.

\textbf{Metode.} Kasus utama adalah sistem metropolitan Jabodetabek. Wilayah Statistik Metropolitan Indonesia digunakan sebagai sabuk pembanding. DKI Jakarta diperlakukan sebagai satu inti dengan batas administrasi tetap. Kecamatan menjadi unit utama jika resolusi data mendukung. Data pada tingkat kabupaten/kota, zona transportasi, raster, dan titik dipertahankan pada resolusi aslinya sebelum dilakukan agregasi yang dapat ditelusuri. Analisis membedakan data yang teramati, tabulasi agregat, hasil pemodelan, dan proksi.

\textbf{Cakupan wilayah dan waktu.} Baseline utama menggunakan tahun 2024. Data tahun lain digunakan sebagai konteks atau pemeriksaan perubahan, dengan tahun dan perbedaan unit ditampilkan secara terbuka. Wilayah mencakup DKI Jakarta, kabupaten/kota dalam kawasan utama Jabodetabek, pusat sekunder yang ditetapkan melalui prosedur yang terdokumentasi, serta sabuk pembanding yang diperlukan untuk menguji fungsi relatif.

\textbf{Jenis data dan variabel.} Penelitian memerlukan data penduduk, luas, kepadatan, integrasi dan perjalanan, fungsi layanan lokal, pekerjaan menurut lokasi tempat kerja dan sektor jika tersedia, aksesibilitas jaringan, serta karakteristik pusat sekunder. Podes digunakan untuk mengukur kehadiran dan akses layanan, bukan sebagai ukuran langsung produktivitas atau jumlah pekerjaan. Data komuter pada tingkat kabupaten/kota tidak akan dipaksa menjadi estimasi kecamatan. Data OD yang dimodelkan akan diberi label sintetis dan diuji terhadap kendala yang tersedia.

\textbf{Rencana hasil.} Hasil penelitian berupa profil integrasi dan fungsi lokal, tipologi kondisi relasional, peta gradien, uji sensitivitas bobot dan ambang, serta penjelajah skenario berbasis WebGL. Penjelajah tersebut menghitung ulang dan menampilkan hasil secara interaktif. Ia bukan sumber data waktu nyata, ramalan kebijakan, atau simulasi perubahan kewenangan pemerintahan.
}

\newcommand{\LampiranAbstraksi}{%
  \clearpage
  \section*{Lampiran 2. Abstraksi penggunaan data}
  \AbstraksiPenelitian
}

\newcommand{\TandaTanganDTAP}{%
  Hormat kami,\\[2.2cm]
  \parbox[t]{0.85\textwidth}{[NAMA PENANDATANGAN]}\\
  \parbox[t]{0.85\textwidth}{[JABATAN PENANDATANGAN SESUAI ALUR PERSURATAN DTAP]}
}

\newcommand{\TandaTanganBPS}{%
  Hormat kami,\\[2.2cm]
  \parbox[t]{0.85\textwidth}{[NAMA PENANDATANGAN]}\\
  \parbox[t]{0.85\textwidth}{[DEKAN FAKULTAS TEKNIK ATAU PEJABAT SETINGKAT YANG DITERIMA BPS]}
}


% Submit through the current PPID Kementerian Perhubungan channel.
% The former BPTJ study name is used as an archive reference, not as the
% assumed current custodian.

\begin{document}
\KopSurat

\begin{flushright}
Yogyakarta, [TANGGAL SURAT]
\end{flushright}

\begin{tabularx}{\textwidth}{@{}P{0.20\textwidth}Y@{}}
  Nomor & [DIISI OLEH UNIT PERSURATAN] \\
  Lampiran & Dua lampiran \\
  Hal & Permohonan informasi dan data agregat asal dan tujuan Kajian Transportasi Jabodetabek 2023
\end{tabularx}

Yth. Pejabat Pengelola Informasi dan Dokumentasi\\
Kementerian Perhubungan Republik Indonesia\\
c.q. [UNIT PENGUASA INFORMASI SESUAI PORTAL PPID]\\
Jakarta

Dengan hormat,

Sehubungan dengan penyusunan Tugas Akhir pada Program Sarjana Perencanaan Wilayah dan Kota, Departemen Teknik Arsitektur dan Perencanaan, Fakultas Teknik, Universitas Gadjah Mada, kami mengajukan permohonan informasi untuk mahasiswa berikut.

\IdentitasMahasiswa

Penelitian berjudul sementara ``\JudulSementara''. Penelitian ini memerlukan bukti tentang hubungan perjalanan antara DKI Jakarta dan pusat-pusat sekunder di kawasan Jabodetabek. Laporan kajian transportasi Jabodetabek 2023 yang tersedia untuk publik menyajikan tabel bangkitan dan tarikan pada 182 kecamatan. Kami ingin memahami data, zona, dan metode yang mendasari tabel tersebut tanpa meminta data perjalanan individu.

Melalui mekanisme layanan informasi publik, kami memohon data agregat dan dokumentasi pada Lampiran 1. Jika arsip tersebut berada pada unit lain atau merupakan keluaran kajian sebelum perubahan organisasi, mohon permohonan ini diteruskan kepada unit yang menguasai informasi atau diberi petunjuk mengenai kanal yang tepat.

\PernyataanPenggunaanData

\PernyataanTanpaBiaya

Jika matriks lengkap tidak dapat diberikan, bentuk pengganti pada Lampiran 1 tetap dapat digunakan untuk penelitian. Kami juga menerima penjelasan tertulis mengenai bagian yang dapat diberikan, bagian yang dikecualikan, dan prosedur keberatan yang berlaku.

Demikian permohonan ini kami sampaikan. Atas perhatian dan bantuan Bapak/Ibu, kami mengucapkan terima kasih.

\vspace{8pt}
Tembusan:
\begin{enumerate}
  \item Dosen pembimbing;
  \item Ketua Program Studi Sarjana Perencanaan Wilayah dan Kota;
  \item Arsip.
\end{enumerate}

\TandaTanganDTAP

\clearpage
\section*{Lampiran 1. Rincian informasi yang dimohon}

\renewcommand{\arraystretch}{1.25}
\begin{tabularx}{\textwidth}{@{}P{0.12\textwidth}P{0.34\textwidth}Y@{}}
\toprule
Prioritas & Informasi atau data & Unit dan kelengkapan yang dimohon \\
\midrule
P0 & Matriks asal dan tujuan lengkap yang mendasari kajian transportasi Jabodetabek 2023 & Diutamakan 182 x 182 kecamatan. Mohon dicantumkan unit arus, periode, hari rujukan, tujuan perjalanan, moda, dan status data teramati atau termodelkan. \\
P0 & Definisi zona dan padanan wilayah & Kode dan nama zona, geometri atau crosswalk zona ke kecamatan, sistem koordinat, tahun batas, serta catatan zona yang memotong atau berubah dari batas administrasi. \\
P0 & Kamus data dan asal-usul matriks & Sumber input, definisi variabel, faktor ekspansi, pembulatan, nilai nol atau hilang, serta aturan penyamaran. \\
P1 & Metode, kalibrasi, dan validasi & Fungsi hambatan, sumber jaringan dan waktu tempuh, target kalibrasi, indikator evaluasi, galat yang diketahui, dan keterbatasan model. \\
P1 & Keluaran agregat sebagai pengganti & Jika matriks lengkap tidak tersedia, mohon pertimbangkan arus antarkelompok pusat, arus menuju dan dari DKI Jakarta, tiga sampai lima pusat sekunder prioritas, atau tabulasi tujuan dominan tanpa data pribadi. \\
\bottomrule
\end{tabularx}

\vspace{8pt}
Permohonan ini tidak mencakup nama, alamat, NIK, nomor kontak, jejak perjalanan individu, data telepon pintar, atau pengenal usaha. Kami tidak meminta data yang berada di bawah pembatasan kerahasiaan pihak ketiga.

\LampiranAbstraksi

\end{document}
